\section{Conclusion}
\label{sec:conclusion}

In this work, we investigated whether self-reflection in LLMs induces structured shifts in internal representations through the lens of \textit{representation drift} and \textit{linear separability}. We found that while representation drift occurs universally across various revision conditions, it is not a direct indicator of cognitive restructuring. For the 1.5B model, autonomous self-critique fails to improve the linear separability of correct reasoning, yielding results comparable to a simple clarity rewrite. In contrast, providing an external critique triggers a significant restructuring of the latent reasoning space, nearly doubling the linear probe AUC for correctness. Paradoxically, the 0.5B model shows a robust autonomous capacity for restructuring its internal states, suggesting that the dynamics of self-reflection are highly model-specific. Our results provide a novel framework for assessing the "depth" of LLM self-reflection and underscore the importance of distinguishing between shallow representational changes and meaningful internal restructuring for building the next generation of recursive self-improving AI systems.
